\documentclass{report}
\usepackage{amsmath}
\usepackage{graphicx}
\graphicspath{{figures/}}
\usepackage{subcaption}
\usepackage{adjustbox}
\usepackage{float}
\usepackage{xcolor}
\usepackage{array}
\usepackage{siunitx}
\usepackage{caption}
\usepackage{threeparttable}
\usepackage[toc,page]{appendix}
\renewcommand{\appendixtocname}{Appendix}  % Changes "Appendices" to "Appendix" in ToC
\renewcommand{\appendixpagename}{Appendix} % Changes "Appendices" to "Appendix" in the document header
\definecolor{linkblue}{RGB}{6,125,233}
\usepackage{hyperref}
\hypersetup{
    colorlinks=true,
    linkcolor=linkblue,
    urlcolor=linkblue,
    citecolor=linkblue
}
\usepackage{booktabs}
\title{Convergence Study of a Propeller}
\author{Ethan Creer}
\date{10 December 2025}
\begin{document}

\maketitle

\begin{abstract}
    This report presents a series of mathematical and pragmatic convergence analysis, the principles of which can be applied to a variety of computational fluid dynamics applications.
    \bigskip
    The Weber Wing convergence study 
    
    
    ADD ABSTRACT FOR WING CONVERGENCE STUDY.

    REVIEW LOGIC OF WEBER WING CHAPTER.

    REWRITE NARRATIVE STYLE OF WEBER WING CHAPTER.

    ADD GLYPH FIGURE OF WEBER WING.

    TRANSITION BETWEEN CHAPTERS 1 AND 2.

    \bigskip

    The second study was on an APC 10$\times$7E propeller via the \texttt{FLOWUnsteady} simulation package. The objective was to find optimal discretization parameters that use an efficient computational cost and allocation of supercomputer resources to determine the converged thrust coefficient. Temporal convergence was established via Richardson extrapolation, with 240 time steps per revolution over 3 revolutions as the converged result. Spatial convergence, constrained by numerical and geometric instability, was limited by a numerical stability ceiling. As such, the number of spanwise elements per blade was limited to \texttt{n} = 46 and further convergence was not possible. Analyzing the number of particles shed per step revealed that limiting \texttt{p\_per\_step} to 3 provided adequate thrust coefficient convergence and effectively reduced computation time. A verification simulation, with increased parameter values, determined that computational time was reduced by 85\%, from approximately 8 hours to 1 hour, with a thrust coefficient difference of only 0.27\%. The significant drop in computational time and cost demonstrates that the suggested parameters provide high-fidelity thrust coefficient convergence for this rotor geometry.
\end{abstract}

\tableofcontents


\chapter{Weber Wing Convergence Study}
\section{Problem Setup}

Before performing a preliminary convergence study analysis, it is necessary to define the wing geometry and simulation parameters used in the study.
These parameters are defined as found in Table~\ref{tab:geometry_weber}

\begin{table}[H]
    \centering
    \caption{The Weber wing geometry and simulation parameters with their given values.}\label{tab:geometry_weber}
    \begin{tabular}{c l l}
        \toprule
        \textbf{Parameter} & \textbf{Meaning} & \textbf{Value} \\
        \midrule
        $\alpha$ & Angle of Attack & 4.2$^\circ$ \\
        $V_{\infty}$ & Freestream Velocity & 49.7 m/s \\
        $\rho$ & Air Density & 0.93 kg/$m^3$ \\
        b & Span Length & 2.489 $m$ \\
        AR & Aspect Ratio b/${c_{tip}}$ & 5.0 \\
        TR & Taper Ratio ${c_{tip}}/{c_{root}}$ & 1.0 \\
        $\lambda$ & Sweep Angle & 45.0$^\circ$ \\
        $\gamma$ & Dihedral & 0.0$^\circ$ \\

    \end{tabular}
\end{table}

With the wing now defined, we are now set to begin the convergence study on the wing. We are primarily interested in the variables listed in Table~\ref{tab:adjustables}.
The convergence mode for each is listed in the same table giving us a preliminary idea of how each variable will affect the convergence of coefficient of lift and drag values.


\begin{table}[H]
    \centering
    \caption{A list of adjustable variables that alter the resolution of the output variables.}\label{tab:adjustables}
    \begin{adjustbox}{center}
    \begin{tabular}{l l c}
        \toprule
        \textbf{Variable} & \textbf{Meaning} & \textbf{Convergence Mode} \\
        \midrule
        \texttt{wakelength} & Length particles travel down the wake & Temporal \\
        \texttt{nsteps} & Number of Time Steps & Temporal \\
        \texttt{n} & Number of Spanwise Elements per Side & Spatial \\
        \texttt{p\_per\_step} & Particles shed per step & Spatial \\
        \bottomrule
    \end{tabular}
    \end{adjustbox}
\end{table}

Our convergence study will follow a systematic approach, gradually iterating through variables to suggest final parameters such that the wing will be appropriately discretized and sufficiently converged.

\section{Wakelength}
The first run through of the program consists of varying the wakelength of the wing. Obtaining convergence via wakelength varying is important as the wing simulation needs to be run long enough for the wake to settle in and the furthest aft particles no longer significantly influence the lift and drag.

In our study, we vary the wakelength as a multiplicative factor of the span length given by Equation~\ref{eq:wakelength}. 
\begin{equation}\label{eq:wakelength}
wakelength = i \times{} b
\end{equation}

We perform a sweep across the $i$ value from 0.75 to 5.0 in 0.25-step increments. However, for computational efficiency, once when each of the coefficients of lift and drag become sufficiently converged, we break from this loop.
Determining the metric used for convergence indication was also an important process. Here, we chose to use the coefficient of variation value for the 5 tail values of $C_L$ and $C_D$ at the end of each simulation.
When each of these metrics fell below a threshold value (in this case 0.002 --- or 0.2\%), the simulation was considered converged, and this wakelength value was stored as a global variable to be referenced in subsequent convergence studies.
In our sweep, this converged value ended up being 3.11125.
The coefficient of variation is essentially a normalized standard deviation by taking the standard deviation of the given subset and dividing by its associated mean. 
The coefficient of variation is helpful in the analysis of multiple datasets, and because the variable sweep naturally produces a variety of different datasets, it is appropriate for this study.

Wakelength convergence is particularly important when it comes to temporal convergence. It is essentially testing for how long the wakelength needs to be for the $C_L$ and $C_D$ values to flatten out at the end of the simulation.
This will be shown near the end of the study with the final proposed parameter values.

\section{Number of Time Steps (\texttt{nsteps})}
The second parameter of interest in our convergence study is that of \texttt{nsteps}. This value dictates the number of time steps that take place from the wakelength being its shortest to its longest length.
Increasing \texttt{nsteps} does not shorten or lengthen the wakelength, but it does increase the resolution of the wakelength, allowing for more accurate $C_L$ and $C_D$ calculations.
For our study, we fix the wakelength to its converged value from the previous step and then sweep \texttt{nsteps} from 10 to 100, this time in increments of 10.
To determine convergence, we use the same metric used in the wakelength convergence test --- that is, we measure the coefficient of variation for the final 5 values of $C_L$ and $C_D$ for each run.
Again, once when this coefficient falls below 0.002 --- or 0.2\%, we consider the study converged, store the variable in the global scope for future reference, and transition to the next test.
In this sweep, the wing coefficients converged at a value of $\texttt{nsteps} = 20$.


\section{Spanwise Elements per Side (\texttt{n})}
The next parameter of interest, \texttt{n}, is an important spatial convergence parameter. In essence, it is a measure of how finely the blade is discretized in spanwise panels.
The greater the value of \texttt{n}, the greater the number of spanwise panels are present per side of the wing.
With the wakelength and \texttt{nsteps} values set from the previous sweeps, we now sweep the value of \texttt{n} from 10 to 100, again in increments of 10.
However, instead of using the coefficient of variation as the metric for convergence, we instead use a run over run percent change calculation.
We do this because of the fundamental goal of this test. With temporal convergence already worked out we are now testing for spatial convergence --- which in this test means the final values of $C_L$ and $C_D$ not changing significantly from run to run.
Instead of using the singular final value for both of these coefficients, we attempt to account for numerical noise by taking the average of the last 3 recorded values and then performing the percent change calculation according to Equation~\ref{eq:percent_change}.
In effect, we take the difference between the current run's coefficient value and the previous run's coefficient value, divide it by the previous run's value, and multiply by 100\%.
The tolerance for this study was set to be 0.25\%, and once the run over run percent change for both coefficient values fell below this tolerance, the for loop broke away, and the values were recorded as global variables for future analysis.
For this study, the converged value for \texttt{n} was 40 --- meaning 40 spanwise panels on each side of the simple wing.

\begin{equation}\label{eq:percent_change}
\% \text{ Change} = \frac{\text{Coeff}_{\text{curr}} - \text{Coeff}_{\text{prev}}}{\text{Coeff}_{\text{prev}}} \times{} 100\%
\end{equation}


\section{Particles Shed per Step (\texttt{p\_per\_step})}
A final parameter of interest necessitates discussion in this convergence study --- that of p\_per\_step.
This parameter did not receive its own sweep, but rather was internally updated across each run of the studies various sweep.
An understanding of what this parameter physically represents can shed light on its importance in this convergence study.
p\_per\_step is best described as the number of particle shed per step, per spanwise panel and is a key parameter in limiting particle overlap on each spanwise panel.
In our study, we set the minimum p\_per\_step value to be the result of Equation~\ref{eq:pperstep}, rounded up to the nearest integer.

\begin{equation}\label{eq:pperstep}
\frac{2 \cdot n \cdot \lambda_{vpm} \cdot V_\infty \cdot ttot}{nsteps \cdot b}
\end{equation}

Here, $\lambda_{vpm}$ is a particle overlap parameter and is fixed for the entire study. $ttot$ is also fixed, coming from the wakelength convergence study.

The overriding principle inherent in this parameter, is that its value must change with each run, such that the true independent variable being swept can be isolated and analyzed.
The constraint here is that particles should not overlap each other (as good as possible), and because p\_per\_step is a function of \texttt{n}, \texttt{nsteps}, and (indirectly) wakelength, its value must be updated with each run.
It's essentially its own dependent variable and a necessary constraint to insure uniformity across the variable sweep.

As such, this variable, was updated with each new variable configuration to provide the most accurate coefficient values.
Though not an explicitly swept parameter in this study, it did require an understanding of its role in each wing discretization.

\section{Weber Wing Recommended Parameter Values and Verification}

With the above convergence studies performed, and an understanding of the meaning of \texttt{p\_per\_step}, we are now ready to present a list of parameter values that would provide adequate coefficient of lift and drag values in a study.
Table~\ref{tab:proposed} presents these values below.



\begin{table}[H]
    \centering
    \caption{A final proposed list of parameter values that can be used with sufficient confidence for producing converged \texorpdfstring{$C_L$}{CL} and \texorpdfstring{$C_D$}{CD} values on a simple wing.}\label{tab:proposed}
    \begin{tabular}{l c}
        \toprule
        \textbf{Variable} & \textbf{Proposed Value} \\
        \midrule
        \texttt{wakelength} & 3.111 \\
        \texttt{nsteps} & 20 \\
        \texttt{n} & 40 \\
        \texttt{p\_per\_step} & 10 \\
        \bottomrule
    \end{tabular}
\end{table}

\section{Verification Tests}
With this proposed set, we now set to gain greater confidence in these converged values by performing a series of verification tests.
These tests involve individually increasing each variable by a certain factor (in our case by 1.5 or 2), re-running the solver, and analyzing whether the values for $C_L$ or $C_D$ changed significantly.


\begin{table}[H]
    \centering
    \caption{Percent change values in \texorpdfstring{$C_L$}{CL} and \texorpdfstring{$C_D$}{CD} when varying an independent variable by a multiplicative factor.}\label{tab:verification}
    \begin{adjustbox}{center}
    \begin{tabular}{c c c}
        \toprule
        \textbf{Variable (Multiplicative Factor)} & \textbf{\% Change in $C_L$} & \textbf{\% Change in $C_D$}\\
        \midrule
        \texttt{nsteps} (2) & 0.3739 & 0.2078 \\
        \texttt{p\_per\_step} (1.5) & 0.0307 & 0.0311 \\
        \texttt{wakelength} (1.5) & 0.4477 & 0.4847 \\
        \bottomrule
    \end{tabular}
    \end{adjustbox}
\end{table}

The data presented in Table~\ref{tab:verification} depicts the percent change in $C_L$ and $C_D$ when independently varying the listed parameters.
With a tolerance of 0.5\%, we can see that each verification test succeeds in returning sufficiently converged values.


\begin{figure}
    \centering
    \includegraphics[width=0.9\textwidth]{CL_vs_time_proposed.png}
    \caption{A visual representation of temporal convergence for the coefficient of lift value. As the wake develops in time, the coefficient level off to become steady.}\label{fig:cl_plot}
\end{figure}

\begin{figure}
    \centering
    \includegraphics[width=0.9\textwidth]{CD_vs_time_proposed.png}
    \caption{A visual representation of temporal convergence for the coefficient of drag value. As the wake develops in time, the coefficient begins to levels off at a steady value. Note the y-axis scale here, which deals with very small numbers.}\label{fig:cd_plot}
\end{figure}

As one final visual confirmation of the converged values, we examine both a $C_L$ vs Time and $C_D$ vs Time plot to confirm temporal convergence for each.
This is depicted in Figures~\ref{fig:cl_plot}~and~\ref{fig:cd_plot}. In the coefficient of lift vs time plot we can see that the coefficient of lift levels off as the wakelength fully develops in time.
As such, we have further confirmation that these parameter settings provide sufficient resolution for convergence for coefficient of lift.

For the coefficient of drag plot, we note some unique behavior towards the beginning of the simulation, with the coefficient of drag value shooting up, and then exponentially decaying in time. This plot does not have as defined of a leveling off for coefficient of drag value, but we note that the y-axis has a very small scale, that varies from $4.8 \times 10^{-3}$ to $4.9 \times 10^{-3}$.
This small range could in part explain why the graph does not appear to level off. Our study relied on coefficient of variation values in addition to run over run percent change in values to determine convergence, and with such small magnitude in numbers, this could look as if it prematurely converged.

If we examine the \texttt{nsteps} verification test and its associated $C_D$ vs Time plot, seen in Figure~\ref{fig:cd_plot_verified}, we note that there is a more defined ``leveling-off'' of the coefficient of drag value.

\begin{figure}
    \centering
    \includegraphics[width=0.9\textwidth]{CD_vs_time_verification1.png}
    \caption{A visual representation of temporal convergence for the coefficient of drag after doubling the \texttt{nsteps} value from its converged proposal. As the wake develops in time, the coefficient begins to levels off at a steady value.}\label{fig:cd_plot_verified}
\end{figure}

Thus, if greater convergence is desired, one can use the doubled \texttt{nsteps} value, which will provide more of a visual level off.
However, for most applications, the proposed values provide sufficient convergence, and will only deviate from true coefficient values by a small amount.

\section{Weber Wing Conclusions}
Ultimately, a convergence study provides crucial analysis for the variety of parameters that come with aerodynamic studies.
Without the convergence test just performed, one could unnecessarily use solver settings that directly lead to a more computationally expensive solve.
Though this does not take significant time in the case of a simple wing, the computational time greatly increases as the wing and/or propeller system becomes more complicated.
Thus, from an optimization standpoint, it is important that the right parameter settings are determined such that with some level of confidence we can treat the resulting $C_L$ and $C_D$ values as true to their respective conditions.


\chapter{10$\times$7E Propeller Convergence Study}
\section{Problem Setup}

To begin our convergence study analysis, it is necessary to define the propeller geometry and simulation parameters used in the study.
These parameters are defined in Table~\ref{tab:geometry_prop}, and a 3D loft of the rotors is visualized in Figure~\ref{fig:lofts}.

\begin{table}[H]
    \centering
    \caption{Propeller geometry parameters with their given values. For simplicity, only the overall rotor is defined and referenced. The blade geometry is dictated by its model (an APC 10$\times$7E rotor) with relevant geometry found on the BYU Flow Lab \href{https://github.com/byuflowlab/FLOWUnsteady/tree/master/database/rotors}{GitHub}.}\label{tab:geometry_prop}
    \begin{tabular}{l l l}
        \toprule
        \textbf{Parameter} & \textbf{Description} & \textbf{Value} \\
        \midrule
        $R_{tip}$ & Radius of Blade Tip (m) & 0.127 \\
        $R_{hub}$ & Radius of Hub (m) & 0.00953 \\
        $B$ & Number of Blades & 2 \\
    \end{tabular}
\end{table}

As will be discussed further in Section~\ref{future}, we discovered a discrepancy between the rotor geometry linked on the manufacturer's website and what is described in the GitHub files. As such, this report will focus on the theoretical convergence using the parameters provided on GitHub. Further research will be needed to compare experimental data with our converged results.

\begin{figure}[H]
    \centering
    \includegraphics[width=0.9\textwidth]{lofts.png}
    \caption{A 3D representation of the rotor blade geometry.}\label{fig:lofts}
\end{figure}


With the rotor geometry defined, we can take the initial steps in our convergence study. We are primarily interested in the variables listed in Table~\ref{tab:adjustables_rotor}.
The convergence mode for each is listed in the same table giving us a preliminary idea of how each variable will affect the convergence of the thrust coefficient.


\begin{table}[H]
    \centering
    \caption{A list of adjustable variables that alter the resolution of the output thrust coefficient.}\label{tab:adjustables_rotor}
    \begin{adjustbox}{center}
    \begin{tabular}{l l c}
        \toprule
        \textbf{Variable} & \textbf{Meaning} & \textbf{Convergence Mode} \\
        \midrule
        \texttt{nrevs} & Number of Propeller Revolutions & Temporal \\
        \texttt{nsteps\_per\_rev} & Number of Time Steps per Propeller Revolution & Temporal \\
        \texttt{n} & Number of Spanwise Elements per Propeller Blade & Spatial \\
        \texttt{p\_per\_step} & Particles shed per step & Spatial \\
        \bottomrule
    \end{tabular}
    \end{adjustbox}
\end{table}

These variables will be the focus of our convergence study, and we will follow a systematic approach, gradually iterating through each. We will eventually suggest final parameters that can be used to sufficiently discretize the propeller blades and converge the thrust coefficient. Whether these parameters work on any blade and its respective geometry is another question that further research will be done on in the future. However, the core processes and principles presented in our report can be adapted and applied to any propeller.

\section{Number of Revolutions (\texttt{nrevs})}
The convergence of propeller revolutions was initially evaluated with a visual inspection of a $C_T$ vs time plot, seen in Figure~\ref{fig:ct_vs_time_init}.
We examine the plot for the point where the thrust coefficient reaches a steady state (or begins to level off). We can then estimate the number of revolutions (to the nearest whole number, as required by the code documentation) before the coefficient levels off in time.

The simulation results depicted in Figure~\ref{fig:ct_vs_time_init} ran for a total of 8 revolutions, and we see that the plot begins to level off about a third of the way into the simulation time. From a mathematical standpoint, Table~\ref{tab:ct_vs_time_percent} presents the revolution on revolution percent change in the thrust coefficient.

\begin{table}[H]
    \centering
    \caption{A decline in revolution on revolution percent change with increasing number of revolutions.}\label{tab:ct_vs_time_percent}
    \begin{adjustbox}{center}
    \begin{tabular}{c c}
        \toprule
        \textbf{Revolution Number} & \textbf{Percent Change (\%)} \\
        \midrule
        1 & N/A \\
        2 &  -9.8827 \\
        3 &  -0.6778 \\
        4 &  -0.2286 \\
        5 &  -0.0974 \\
        6 &  -0.0564 \\
        7 &  -0.0345 \\
        8 &  -0.0134 \\
        \bottomrule
    \end{tabular}
    \end{adjustbox}
\end{table}

Noting that revolution 3 is accompanied by only a 0.6778\% change in thrust coefficient, we will take that to be our converged parameter value.
Though this parameter study was done on relatively low discretization settings, the value will be confirmed again in the verification study of our converged recommended values presented in Section~\ref{Verification}.

\begin{figure}[H]
    \centering
    \includegraphics[width=0.9\textwidth]{CT_vs_time_nrevs8.png}
    \caption{A plot representing the temporal convergence of the thrust coefficient. Convergence depicted by the plot leveling off.}\label{fig:ct_vs_time_init}
\end{figure}


\section[Number of Time Steps per Revolution]{Number of Time Steps per Revolution \\ \normalfont(\texttt{nsteps\_per\_rev}, hereafter referenced as \texttt{nsteps})}
With the number of revolutions sufficiently converged, we finalize our temporal convergence with a study on the number of time steps per revolution.
This parameter directly influences the resolution of the simulation and has a strong effect on the final coefficient values. It essentially defines the degrees each step the rotor blades traverse.

Our approach to converging this parameter involves leveraging the principle of Richardson extrapolation. This technique requires at least 2 data points from preliminary runs and increases in accuracy with the more data points supplied. 
Its purpose is to extrapolate beyond a given dataset's range and provide an estimated value for the dependent variable (in this case the thrust coefficient) as the independent variable approaches its continuum limit (in this case, as $\Delta t \rightarrow 0$). Additionally, this extrapolation theory is based on the assumption that changing the time step between runs introduces a different error from the true value. Equations~\ref{eq:richardson1} and~\ref{eq:richardson2} begin to illustrate this concept:

\begin{equation}\label{eq:richardson1}
f_{\Delta t} = f_{exact} + C(\Delta t)^p
\end{equation}
\begin{equation}\label{eq:richardson2}
f_{\Delta t/2} = f_{\text{exact}} + C\left(\frac{\Delta t}{2}\right)^{p}
\end{equation}

where $f_{\Delta t}$ is the estimated value from the simulation run, $f_{exact}$ is the true value, $C$ is some constant, and $p$ is the order of accuracy from the numerical method. Because we used more than 3 data points, an observed order of accuracy can be calculated through the additional $f_{\Delta t}$ equations generated.
These two equations can be algebraically manipulated to eliminate C and solve for $f_{exact}$. Essentially, it takes advantage of the fact that a smaller time step results in a smaller error to extrapolate for the true value.

Once the Richardson extrapolation is complete, the extrapolated $C_T$ value can then be used to estimate what value of \texttt{nsteps} would give us a sufficiently converged thrust coefficient value.

This Richardson extrapolation process is visualized for our case in Figure~\ref{fig:nsteps_per_rev_richardson}, where the red dashed line represents the Richardson extrapolated $C_T$ value as \texttt{nsteps} approaches $\infty$.
\begin{figure}
    \centering
    \includegraphics[width=0.9\textwidth]{FINAL320_nsteps_richardson.png}
    \caption{Richardson's extrapolation of $C_T$ as \texttt{nsteps} $\rightarrow \infty$.}\label{fig:nsteps_per_rev_richardson}
\end{figure}

For this study, we used the following values for \texttt{nsteps}: 20, 40, 80, 160, and 320. Using the extrapolated value of 0.087177, we can tabulate the data from each run's final $C_T$ value to gain insight into the converged value of \texttt{nsteps}.
This tabulated data is seen in Table~\ref{tab:nsteps_per_rev_richardson} where the converged $C_T$ value is taken as an average of the final 5 values of $C_T$. This is done to account for the slight oscillations present in the data near the end of the simulation.

\begin{table}[H]
    \centering
    \caption{Tabulated data from Richardson extrapolation of $C_T$ with varying \texttt{nsteps}.}\label{tab:nsteps_per_rev_richardson}
    \begin{adjustbox}{center}
    \begin{tabular}{c c c}
        \toprule
        \textbf{\texttt{nsteps}} & \textbf{Time Averaged $C_T$} & \textbf{Percent Error from Extrapolated $C_T$ (\%)} \\
        \midrule
        20 & 0.08117 & 6.8931 \\
        40 & 0.08273 & 5.0965 \\
        80 & 0.08417 & 3.4537 \\
        160 & 0.08549 & 1.9303 \\
        320 & 0.08629 & 1.0132 \\
        \bottomrule
    \end{tabular}
    \end{adjustbox}
\end{table}

For this parameter, we set a target percent error of 1.5\% and use simple interpolation to yield a converged value for the \texttt{nsteps} parameter of approximately 235.
For simplicity, we round up to the nearest 10, and settle on a converged value of 240 for \texttt{nsteps}. This too will be verified later in our verification run.


\section{Number of Spanwise Elements per Propeller Blade (\texttt{n})}\label{n_convergence}
With the parameter settings for temporal convergence known, we then transitioned to spatial convergence by first addressing the number of spanwise elements per blade, \texttt{n}.
Knowing that the aerodynamic forces vary across the span of the propeller, it is important to capture the varying particle vorticity with a sufficiently discretized blade.
The process to adequately capture this distributed particle vorticity involved significant debugging, of which more information can be found in Appendix~\ref{app:debugging}.

\bigskip

After the debugging process, converging the number of spanwise panels per blade can occur.
This is done similarly to the manner in which \texttt{nsteps} was converged. The script was configured to reference the converged \texttt{nsteps} value of 240, and used a job array to run multiple iterations on different supercomputer nodes.

In reference to the Richardson extrapolation, this portion of the study had complications in terms of rotor generation from the \texttt{FLOWVLM} functions and numerical stability in the \texttt{FLOWUnsteady} library.
To counteract these issues, fine-tuning parameters such as the geometric expansion factor, \texttt{r}, the Rotor on VPM (vortex particle method) smoothing radius, \texttt{sigma\_rotor\_surf}, and the VLM (vortex lattice method) relaxation, \texttt{vlm\_rlx}, were implemented.

Further, an RPM ramp up scheme was imposed on the propeller to replace the instantaneous jump from 0 to 9200 RPM\@. This was done to mitigate the effects of the initial particle sheds that were difficult for the simulation to model.
These particles typically had much stronger vortex strengths compared to particles shed later in the run. We also attempted to mitigate these adverse effects by adjusting the \texttt{vlm\_init} argument to be true in the \texttt{run\_simulation ()} function call.

These fine-tuning adjustments required significant computational time and yielded inconclusive results in some instances. However, after a significant number of iterations and careful analysis, we propose that \texttt{n=46} is a strong candidate for thrust coefficient convergence.
This decision involved the careful analysis of Figure~\ref{fig:defense_of_n46} and Table~\ref{tab:n_richardson_percent_error} described below.

\begin{figure}
    \centering
    \includegraphics[width=0.9\textwidth]{defense_of_n46.png}
    \caption{$C_T$ simulations with varying number of spanwise panels. A ±1.5\% convergence band is overlayed to aid in convergence analysis. Purple data points, representing the simulations with adjusted spline parameters are also present illustrating the induced geometric error and numerical stability ceiling.}\label{fig:defense_of_n46}
\end{figure}


\begin{table}[H]
    \centering
    \caption{Tabulated data from the spatial convergence study of $C_T$ by varying \texttt{n}.}\label{tab:n_richardson_percent_error}
    \begin{threeparttable}
        % Using tabular* with \textwidth and \extracolsep{\fill} pushes columns apart to fill the page
        \begin{tabular*}{\textwidth}{@{\extracolsep{\fill}} c c c}
            \toprule
            \textbf{\texttt{n}} & \textbf{Time Averaged $C_T$} & \textbf{Percent Error from n=46 (\%)} \\
            \midrule
            30 & 0.08498 & 1.2263 \\
            32 & 0.08509 & 1.3572 \\
            36 & 0.08526 & 1.5633 \\
            38 & 0.08528 & 1.5817 \\
            40 & 0.08503 & 1.2942 \\
            42 & 0.08444 & 0.5825 \\
            44 & 0.08411 & 0.1966 \\
            46 & 0.08395 & (Reference) \\
            48\tnote{a} & 0.08374 & 0.2515 \\
            50\tnote{a} & 0.08412 & 0.2086 \\
            55\tnote{a} & 0.08593 & 2.3559 \\
            60\tnote{a} & 0.08552 & 1.8716 \\
            65\tnote{a} & 0.08475 & 0.9523 \\
            70\tnote{b} & 0.08449 & 0.6468 \\
            75\tnote{b} & 0.08429 & 0.4119 \\
            80\tnote{b} & 0.08392 & 0.039 \\

            \bottomrule
        \end{tabular*}
        \begin{tablenotes}
            \small
            \item [a] Generated using Cubic Spline ($k=3$) with relaxed smoothing ($s=0.01$) to prevent solver failure
            \item [b] Generated using Cubic Spline ($k=3$) with loose smoothing ($s=0.02$)
        \end{tablenotes}
    \end{threeparttable}
\end{table}

Examining the superscripts of Table~\ref{tab:n_richardson_percent_error}, it should be noted that the VLM package had trouble defining a mathematical spline to generate the blade rotors for high values of \texttt{n}.
The default spline parameters for the \texttt{generate\_rotors()} function were to use a 5th order spline and a smoothing factor of $s = 0.0001$. For highly discretized blades, these strict constraints, were incompatible with the mathematical spline solver. As such, we needed to relax the parameters to allow the solver to successfully define the blade geometry.
This, of course, introduced error into the data because some true rotor geometry was lost by the smoothing factor. This can partially explain the atypical behavior present in Figure~\ref{fig:defense_of_n46}.
The lack of exponential decay (or any macroscopic trend) in the data rendered a Richardson extrapolation effectively useless.
Because the experimental $C_T$ values are scattered throughout the green region of Figure~\ref{fig:defense_of_n46}, careful analysis is needed to describe what value of \texttt{n} should be considered ``converged.''

The original convergence criterion was defined by running a verification run --- multiplying the converged value by 1.5 to prove convergence.
However, knowing the problems described earlier that the code experienced with high values of \texttt{n}, the multiplicative factor would not be possible without introducing geometric and thrust coefficient errors.
Because of the spline smoothing factor constraints imposed by the VLM software, the maximum discretization we can achieve is using $\texttt{n} = 46$ as the converged value. Though this may not represent a mathematically converged value, it does represent the highest fidelity possible due to the numerical stability ceiling.
Because of this upper limit on the value for \texttt{n}, of necessity, the multiplicative factor on this parameter will be dropped in the verification study, noting the limitations of the software.


\section{Particles Shed per Step (\texttt{p\_per\_step})}
An important note should be added on how the \texttt{p\_per\_step}, particles shed per step, parameter was handled throughout this convergence study.
Equation~\ref{eq:p_per_step} represents the mathematical description of \texttt{p\_per\_step}:


\begin{equation}\label{eq:p_per_step}
\texttt{p\_per\_step} = \frac{2 \pi \texttt{n} (\lambda_{vpm})}{\texttt{nsteps}}
% p_per_step = Int(ceil(lambda_vpm * 2 * pi * R / (nsteps_per_rev * R/n
\end{equation}

Here, the $\frac{2 \pi}{\texttt{nsteps}}$ expression in the equation represents the angle the propeller traverses in one timestep.
$\lambda_{vpm}$ represents a particle overlap condition. Essentially the particles must overlap to some degree such that an entire sheet of vorticity can mathematically represent the wake in the software.
In our convergence study, $\lambda_{vpm}$ was set to 2.125, which roughly equates to the particle edges touching the center of their adjacent particles.
Finally, \texttt{n} represents the number of spanwise elements per blade, as discussed in Section~\ref{n_convergence}.

While Equation~\ref{eq:p_per_step} calculates the ideal value for \texttt{p\_per\_step}, to prevent a computationally expensive number of particles shed per step, a limit of 3 was put on the parameter.
For the verification run, discussed in Section~\ref{Verification}, the cap was raised to 6 to prove a converged value. 

\section{Propeller Recommended Parameter Values and Verification}\label{Verification}
With each of the mentioned parameters sufficiently converged, we present in Table~\ref{tab:suggested} the recommended parameter values that sufficiently converge the value for $C_T$.
This takes into account the VLM package limitations described in Section~\ref{n_convergence}. It represents a sufficient simulation fidelity using an efficient allocation of resources while also prioritizing computation time.
\begin{table}[H]
    \centering
    \caption{A list of recommended parameter values to produce a converged simulation along with the values used to verify said parameter values.}\label{tab:suggested}
    \begin{adjustbox}{center}
    \begin{tabular}{l c c}
        \toprule
        \textbf{Parameter} & \textbf{Suggested Value} & \textbf{Verification Value} \\
        \midrule
        \texttt{nrevs} & 3 &  5 \\
        \texttt{nsteps} & 240 & 360 \\
        \texttt{n} & 46  & 46 \\
        \texttt{p\_per\_step} Max & 3 & 6 \\
        \bottomrule
    \end{tabular}
    \end{adjustbox}
\end{table}

Running the simulation with the verification parameters took significant time (approximately 8 hours while allocating 128 cores and 256 GB of RAM).
When running with our suggested parameters, the simulation took only about an hour to run (with only 64 cores and 64 GB of memory).
Thus, the proposed converged parameters yielded an 85\% reduction in computation time while only a 0.27\% change in the thrust coefficient as seen in Table~\ref{tab:verified_stats}.

\begin{table}[H]
    \centering
    \caption{A side by side comparison of the thrust coefficient obtained via the recommended and verification parameters.}\label{tab:verified_stats}
    \begin{adjustbox}{center}
    \begin{tabular}{l c c}
        \toprule
        \textbf{Parameters Used} & \textbf{Thrust Coefficient Obtained} & \textbf{Percent Error (\%)} \\
        \midrule
        Suggested & 0.083948 & 0.2732 \\
        Verification & 0.084177 & (Reference) \\
        \bottomrule
    \end{tabular}
    \end{adjustbox}
\end{table}

This thrust coefficient difference of merely 0.2732\% more than justifies the computation time reduction of nearly 7 hours.
As a final visual representation of our convergence study, Figure~\ref{fig:verified_color_adjust_glyphs} should be referenced. The circulation strength (Gamma Magnitude) color scale reveals that the stronger particle interactions take place at the tips of the rotor blades. This helps to justify the geometric expansion factor of $\texttt{r} = \frac{1}{5}$. In effect, this sets the particle spacing near the tips to be one fifth that of the particle spacing near the hub. More particles near the rotor tips allowed the simulation to capture more of the wake data near the complex interactions that take place there.

\begin{figure}[H]
    \centering
    \includegraphics[width=0.9\textwidth]{n46_glyphs.png}
    \caption{A paraview depiction of the rotor wake for the recommended convergence parameters. Color scale is based on particles closest to the blade.}\label{fig:verified_color_adjust_glyphs}
\end{figure}


\chapter{Future Applications}\label{future}
Going forward, there are several questions yet to be answered. Further research comparing the simulation data to experimental data on the physical propeller is of key interest. However, further research resolving the geometry discrepancy from earlier will be needed before experimental validation can occur.

It will also be interesting to look more into if advance ratio, \texttt{J}, affects the parameter values we suggest for thrust coefficient convergence.
Preliminary findings revealed that adjusting the advance ratio led to a divergence from the experimental results. Again though, we discovered that the experimental data may not necessarily correspond to the actual rotor used in the simulation.
As such, more research would be needed.

Future work will also examine the convergence of the wake's average velocity field. This will aid in the analysis of multiple propellers (and larger wind turbines) lined in series and will lead to fascinating insights.

The convergence study of wind turbines is another goal we have. The same principles from this study can be used on the wind turbines. This will be important as the wind turbines have much larger diameters and as such will be accompanied by a significant increase in computation time and data generation. This makes finding the most efficient parameters all the more important. If we adjust solver parameters without a significant decrease in simulation fidelity, we will be able to save significant computation time and as such be able to run more complex simulations.



\begin{appendices}
\chapter{Computational Stability and Dual Language Management}\label{app:debugging}
Core to the \texttt{FLOWUnsteady} library, is the interplay between multiple coding languages. In particular, the primary numerical solver was executed in a Julia environment, while the background visualization plotting was executed in Python via the \texttt{PyCall} package.
Because Python's memory management system is incompatible with threaded runs, we encountered frequent segmentation faults and race conditions associated with the non-serial execution on supercomputer nodes.
Single threaded runs significantly increased the runtime for the simulation, and so, it was necessary to find a solution that addressed segmentation faults and race conditions while still allowing for parallel processing.
This was accomplished through manual control of memory management in the \texttt{FLOWUnsteady} library. 

A \texttt{ReentrantLock} was implemented on Python function calls to address thread safety, preventing race conditions from occuring while the program ran. It effectively serialized Python plot generation calls in the primarily parallel processing based script.
Further, effective memory management was accomplished by periodically calling the Garbage Collector from the Julia base code. This was done by using the \texttt{GC.gc()} command at the end of looped commands to ensure that memory leaks did not occur.
\end{appendices}
\end{document}